\section*{الفصل \ref{ch:b3:plant-molecular-physiology} --- فيزيولوجيا النبات الجزيئية واستجابات الإجهاد}

\begin{solution}{exo:b3:plant-molecular-physiology:1}
الفيتوكرومات، الأحمر \qty{660}{nm} والأحمر البعيد \qty{730}{nm}:
فك \omterm{def:b3:plant-molecular-physiology:photoreceptors}{الاصفرار}، \omterm{prop:b3:plant-molecular-physiology:photoequilibrium}{وتجنب الظل}، والإنبات. والكريبتوكرومات، الأزرق
\qty{450}{nm}: فك \omterm{def:b3:plant-molecular-physiology:photoreceptors}{الاصفرار}، والساعة، والإزهار. والفوتوتروبينات،
الأزرق: الانحناء نحو الضوء، وفتح الثغور، وحركة الصانعات
اليخضورية. واليخضور يقيس الضوء طاقةً ولا يفعل شيئًا
بالمعلومة؛ وأما المستقبِلات الضوئية فتقرأ اللون والاتجاه
والمدة عند دفوق أدنى بمليون مرة وتغيّر التعبير المورّثي ---
وهو الفرق بين أكل الضوء وقراءته.
\end{solution}

\begin{solution}{exo:b3:plant-molecular-physiology:2}
كل \omterm{def:b3:endocrinology:hormone}{هرمون} يعزز يوبقنة كابت وهدمه: \omterm{def:b3:plant-molecular-physiology:hormones}{فالأوكسين} يلصق \omterm{def:b3:plant-molecular-physiology:hormones}{TIR1}
بالبروتينات Aux/IAA، \omterm{def:b3:plant-molecular-physiology:hormones}{والجبريلين} المرتبط بالمستقبل GID1 يسلّم بروتينات
\omterm{def:b3:plant-molecular-physiology:hormones}{DELLA} إلى ليغاز F-box، والجاسمونات يلصق COI1 بالبروتينات JAZ.
والكابتات جالسة أصلًا على مورّثاتها الهدف والمنشِّطات
متأهبة بجانبها، فلا حاجة إلى صنع بروتين جديد: فهدم الكابت،
وهو ما يفعله البروتيازوم في دقائق، يشغّل المورّثات فورًا.
\end{solution}

\begin{solution}{exo:b3:plant-molecular-physiology:3}
يرتبط ABA بالمستقبل PYR/PYL؛ ويثبّط المعقد الفسفاتاز PP2C؛ فيصير
الكيناز SnRK2 (وهو OST1)، إذ لم يعد يُنزَع فسفاته، فعالًا
ويفسفر قناة الأنيون SLAC1؛ فتغادر الأنيونات، ويزول استقطاب
الغشاء، وتنفتح قنوات البوتاسيوم الخارجة، فيغادر K$^{+}$ ثم
الماء، ويهبط الامتلاء وتترهل الخليتان الحارستان معًا. وأي
حلقة يمكن أن تنكسر: فلا مستقبلات، أو فسفاتاز لا يمكن تثبيطه
(وهو الطافر الكلاسيكي \emph{abi1})، أو لا OST1 أو لا SLAC1 ---
وكل منها يعطي نباتًا ذاويًا تبقى ثغوره مفتوحة في الجفاف.
\end{solution}

\begin{solution}{exo:b3:plant-molecular-physiology:4}
تتعرف PTI جزيئات مشتركة بين أصناف ميكروبية كاملة
(الفلاجيلين، والكيتين) بكينازات مستقبِلة عند سطح الخلية
وتنتهي بجدار مقوًّى، وثغور مغلقة، وتعبير عن مورّثات الدفاع،
ونمو أبطأ، بلا موت خلوي عادةً. وتتعرف ETI بروتينًا مؤثِّرًا
بعينه أو أذاه \omterm{def:b3:plant-molecular-physiology:immunity}{بمستقبل NLR} داخل خلوي وتنتهي بالموت فرط الحساسي
للخلايا المصابة وبنذير جهازي.
\end{solution}

\begin{solution}{exo:b3:plant-molecular-physiology:5}
ضوء النهار $0.87\times 1.15/1.75 = 0.57$؛ وورقة واحدة $0.87\times 0.4/1.0 =
0.35$؛ وظُلّة كثيفة $0.87\times 0.1/0.7 = 0.12$؛ ومصباح متوهج
$0.87\times 0.7/1.3 = 0.47$. ومصباح الفتيلة غني بالأحمر البعيد،
فتكون $\varphi$ دون قيمتها في ضوء النهار وتقرأ البادرة ظلًّا
جزئيًا: فتستطيل. والمصابيح البيضاء الباردة والمصابيح
المشعّة للضوء بلا أحمر بعيد تفعل العكس.
\end{solution}

\begin{solution}{exo:b3:plant-molecular-physiology:6}
الجدار $5.0$، والعصارة الخلوية $7.5$: $(1 + 10^{2.75})/(1 + 10^{0.25}) = 563/2.78
= 203$. والجدار $4.5$: $563/(1 + 10^{-0.25}) = 563/1.56 = 360$ ---
فتحميض الجدار يرفع النسبة المتعادلة خارجًا ويرفع الالتقاط.
والعصارة الخلوية $6.5$ مع جدار $5.5$: $(1 + 10^{1.75})/6.62 = 57/6.62 = 8.6$:
فيهبط الاحتباس خمسة أضعاف، ويتسرب \omterm{def:b3:plant-molecular-physiology:hormones}{الأوكسين} من كل وجه،
وتتفلطح التدرجات القطبية --- وهو أحد أسباب فقدان الجذور
الناقصة الأكسجين اتجاهها.
\end{solution}

\begin{solution}{exo:b3:plant-molecular-physiology:7}
\qty{1}{cm/h} $= \qty{10000}{\micro m/h}$: أي $200$ خلية في الساعة،
و\qty{18}{s} في كل واحدة. والانتشار عبر \qty{50}{\micro m}: $t =
(5\times 10^{-5})^{2}/(2\times 5\times 10^{-10}) = \qty{2.5}{s}$.
فداخل الخلية يُعبَر في ثوانٍ؛ ومعظم المدة \qty{18}{s} يُقضى في
انتظار الحمل إلى الخارج عبر حوامل PIN عند الغشاء القاعدي،
وهي الخطوة المحدِّدة للمعدل.
\end{solution}

\begin{solution}{exo:b3:plant-molecular-physiology:8}
$E/A = 1.6\,\Delta w/\Delta c = 1.6\times\num{16000}/150 = 171$ جزيء ماء
لكل CO$_{2}$. وعند C$_{4}$، $\Delta c = 250$: $102$. وفي يوم جاف،
$\Delta w = 24$: $256$.
\end{solution}

\begin{solution}{exo:b3:plant-molecular-physiology:9}
التجمد يجفف الخلايا بسحب الماء إلى جليد خارج خلوي، فالبرد
والجفاف كلاهما إجهاد تجفاف. وهما يتشاركان رسل ثوانٍ (ABA،
وذرى الكالسيوم، والأكسجين الفعال)، وعوامل استنساخ (فعائلة
CBF/DREB يستحثها الاثنان)، ونواتج: فالديهيدرينات وبروتينات
LEA تحمي الأغشية والبروتينات من فقد الماء، والمذابات
المتوافقة (البرولين، والسكريات) تحفظ الماء، ومضادات الأكسدة.
والنبات الذي صنعها لإجهاد واحد يملكها للآخر.
\end{solution}

\begin{solution}{exo:b3:plant-molecular-physiology:10}
عند \qty{1}{\%} من الشمس الكاملة، $k_{1} + k_{2} = \qty{0.005}{s^{-1}}$: أي
95~\% من التوازن بعد $3/k = \qty{600}{s}$، أي عشر دقائق، في
مقابل ست ثوانٍ عند الظهر. وثماني ساعات أربعة أعمار نصف:
فيبقى $1/16 \approx
\qty{6}{\%}$ من Pfr. والومضة الوجيزة تضبط
$\varphi$ ثم يخمد Pfr، وأما النهار فيبقيه مرتفعًا ساعات؛
والاستجابات التي تقتضي أن يعمل Pfr ساعات (الإنبات، وفك
\omterm{def:b3:plant-molecular-physiology:photoreceptors}{الاصفرار}) تكامل التعرض، فلا تستجيب بذرة قلبها المحراث لحظةً
كما تستجيب بذرة راقدة على السطح.
\end{solution}

\begin{solution}{exo:b3:plant-molecular-physiology:11}
عشرة مواضع في اليوم عند $100$ خلية لكل موضع: أي \num{1000}
خلية، و$10^{-4}$ من الورقة؛ وعلى مدى موسم جزء من في المئة، في
مقابل خسارة الورقة لو انتشرت إصابة بلا كبح. فالخلايا الميتة
لا تكلف شيئًا تقريبًا وتأخذ معها طعام الممرض --- بالنسبة إلى
متعيش الحي، الذي يحتاج إلى خلايا حية. وأما متعيش الميت
فيقتات النسيج الميت: \omterm{def:b3:plant-molecular-physiology:immunity}{فاستجابة فرط الحساسية} تناوله وجبة
وموطئ قدم، وبعضها (\emph{Botrytis}، و\emph{Sclerotinia}) يفرز
ذيفانات تستفزها عمدًا؛ وفي مواجهتها يعتمد النبات على دفاعات
يتوسطها الجاسمونات لا على الموت الخلوي.
\end{solution}

\begin{solution}{exo:b3:plant-molecular-physiology:12}
الإنتاج بمعدل $p$ من الساعة 12 إلى 16. وفي النهار الطويل،
والضوء طوال ذلك، $k = \ln 2/4 = \qty{0.17}{h^{-1}}$: فالمستوى عند
16 هو $(p/k)(1 -
\mathrm{e}^{-4k}) = (p/0.17)(0.5) = 2.9p$. وفي النهار القصير، والظلام
من الساعة 10، $k = \qty{1.39}{h^{-1}}$: $(p/1.39)(1 - \mathrm{e}^{-5.5})
\approx 0.72p$. أي CONSTANS أكثر بأربعة أضعاف في النهار الطويل.
وعتبة بينهما --- ولتكن $2p$، لازمة لتفعيل \emph{FT} --- لا
تُعبَر إلا حين توافق نافذةُ الإنتاج عند الساعة الضوءَ:
فتوافق إيقاع داخلي مع النهار الخارجي يحوّل طول نهار متدرجًا
إلى قرار بالإزهار من نوع الكل أو لا شيء.
\end{solution}

\begin{solution}{pb:b3:plant-molecular-physiology:1}
\textbf{1.} ضوء النهار $0.87\times 1.15/1.75 = 0.57$؛ والظُلّة $0.87\times
0.2/0.8 = 0.22$.
\textbf{2.} $r = 2(1 - \varphi)$: أي \qty{0.86}{mm/h} في ضوء النهار،
و\qty{1.57}{mm/h} في الظل؛ وعلى مدى \qty{48}{h}: $41$ في مقابل \qty{75}{mm}،
أي \qty{34}{mm} إضافية.
\textbf{3.} يهبط الأحمر إلى $0.1\times 1.15 = 0.115$ من الأحمر البعيد،
الذي يهبط إلى $0.8$: فتكون $\zeta = 0.144$، و$\varphi = 0.87\times 0.144/0.744 =
0.17$. نعم: فورقة واحدة تهبط بالنسبة $\varphi$ من $0.57$ إلى $0.17$، أي عميقًا في
مدى \omterm{prop:b3:plant-molecular-physiology:photoequilibrium}{تجنب الظل}.
\textbf{4.} كلا معدلي التحول الضوئي يتدرج مع الدفق، فنسبتهما
لا تتدرج. وتحت الغيوم يكاد الطيف لا يتغير وتبقى
$\varphi$ قرب $0.57$: فالبادرة في العراء في يوم كئيب لا
تستطيل --- إذ تقرأ الجيران، لا السطوع.
\textbf{5.} \qty{95}{\%} بعد $3/(k_{1} + k_{2})$: أي \qty{6}{s} في الشمس
الكاملة، و\qty{600}{s} عند \qty{1}{\%}.
\textbf{6.} \qty{8}{h} $=$ أربعة أعمار نصف، $1/16 = \qty{6.3}{\%}$؛
و\qty{14}{h} $=$ سبعة، $1/128 = \qty{0.8}{\%}$. وقد يخبر Pfr الباقي
عند الفجر بطول الليل، غير أن الارتداد تفاعل حراري يجري أسرع
في الليالي الدافئة، وقيمة البداية تتوقف على ضوء النهار:
فالنباتات تقيس طول الليل بالساعة بدلًا من ذلك.
\textbf{7.} الجدار: $1/(1 + 10^{0.75}) = 0.15$؛ والعصارة الخلوية: $1/(1 +
10^{2.45}) = 0.0035$.
\textbf{8.} $(1 + 282)/(1 + 5.6) = 283/6.6 = 43$. والعصارة الخلوية
$43\times 0.1 = \qty{4.3}{\micro M}$.
\textbf{9.} $\qty{1}{cm/h}/\qty{100}{\micro m} = 100$ خلية في الساعة،
و\qty{36}{s} في كل واحدة.
\textbf{10.} \qty{5}{h}. ويبدأ الانحناء في غضون ساعة لأن إعادة
التوزع المهمة جانبية، عبر ميليمتر من القمة، والاستجابة
موضعية؛ والتيار البعيد المدى يحدد الإمداد، لا التوقيت.
\textbf{11.} الخاصرة السفلية: $60/50 = 1.2$، أي \omterm{def:b3:plant-molecular-physiology:hormones}{أوكسين} أعلى من
المعتاد \qty{20}{\%}، والاستطالة $0.8$؛ والخاصرة العليا: $0.8$ من
المعتاد، والاستطالة $1.2$. فتفوق العليا السفلية نموًا: فينحني
الجذر إلى أسفل.
\textbf{12.} خلايا الفرع دون أمثلها الأوكسيني، \omterm{def:b3:plant-molecular-physiology:hormones}{فالأوكسين} الزائد
في الخاصرة السفلية يعزز نموها: فتفوق السفلية العليا نموًا
وينحني الفرع إلى أعلى.
\textbf{13.} $E = 0.3\times 16 = \qty{4.8}{mmol\,m^{-2}\,s^{-1}}$.
\textbf{14.} $A = (0.3/1.6)\times 150 = \qty{28}{\micro mol\,m^{-2}\,s^{-1}}$؛
و$E/A = 4800/28 = 171$.
\textbf{15.} المساحة \qty{0.002}{m^2}، و\num{43200}~ثانية: فالماء $4.8\times
10^{-3}\times 0.002\times\num{43200} = \qty{0.41}{mol} = \qty{7.5}{g}$؛
وCO$_{2}$ $28\times 10^{-6}\times 0.002\times\num{43200} =
\qty{2.4e-3}{mol}$، أي $2.4\times 10^{-3}/6\times 180 =
\qty{0.073}{g}$ من الغلوكوز --- أي مئة غرام من الماء لغرام من
السكر.
\textbf{16.} $E = 0.05\times 16 = \qty{0.8}{mmol\,m^{-2}\,s^{-1}}$،
و$A = \qty{4.7}{\micro mol\,m^{-2}\,s^{-1}}$، والنسبة لا تزال $171$.
فقد وفّر النبات ستة أسباع مائه وتخلى عن ستة أسباع كربونه:
فإغلاق الصمّام يغيّر المعدل، لا الثمن.
\textbf{17.} عند C$_{4}$: $\Delta c = 250$، و$E/A = 4800/46.9 = 102$. وCAM
ليلًا: $E = 0.3\times 5 = \qty{1.5}{mmol}$، و$A = 28$: أي $E/A = 53$.
\textbf{18.} $E/A = 1.6\,\Delta w/\Delta c$: فالناقلية تُختصر لأن
الغازين يمران بالثقب نفسه. ويقدر النبات على تغيير التدرجات
--- بتركيز CO$_{2}$ داخلًا (C$_{4}$)، وبالانفتاح حين يكون
الهواء رطبًا (CAM، والفجر)، وبتبريد الورقة، وبتغليظ الطبقة
الحدية بالشعيرات --- لا على تغيير فيزياء غازين يتشاركان
ثقبًا واحدًا.
\textbf{19.} \num{1000} خلية في اليوم، أي $10^{-4}$ من الورقة.
\textbf{20.} \num{60000} خلية في 60 يومًا، أي \qty{0.6}{\%}. وإفلات
إصابة واحدة في اليوم تتلف \qty{5}{\%} كان سيستهلك الورقة في
عشرين يومًا: \omterm{def:b3:plant-molecular-physiology:immunity}{فاستجابة فرط الحساسية} تكلف جزءًا من مئة مما
تكلفه إصابة واحدة مفلتة.
\textbf{21.} $0.2\times 5 = \qty{1}{\%}$ من التركيب الضوئي. ولو
أُبقيت مشغَّلة دائمًا لكلفت \qty{5}{\%} ولفاق النبات جيرانٌ
ينفقونها على الأوراق؛ فالدفاع المستحَث لا يؤتي ثمنه إلا حين
يكون التهديد حاضرًا.
\textbf{22.} يكسب الاختفاء عن ذلك المستقبل ويخمج الصنف
المقاوم؛ ويخسر ما كان المؤثِّر يفعله (كبت PTI)، فيكون أضعف
على النباتات التي لا \omterm{def:b3:plant-molecular-physiology:immunity}{مستقبل NLR} لديها. وتصير جمهرة النبات
تحابي المستقبلات ضد المؤثِّرات الباقية، وتتراجع مورّثة
المقاومة القديمة، وقد صارت عديمة الفائدة: أي دوران معتمد على
التواتر في مورّثات الطرفين.
\textbf{23.} المحاكي (الكورونَتين) يفعّل سبيل الجاسمونات، الذي
يعاكس سبيل الساليسيلات؛ والساليسيلات هي الدفاع ضد متعيشات
الحي كالبكتيريا نفسها، والمحاكي يعيد أيضًا فتح الثغور التي
كان النبات قد أغلقها. فالممرض يقلب ذراعًا من الجهاز المناعي
على الأخرى.
\textbf{24.} مورّثة مقاومة واحدة على ملايين النباتات المتطابقة
انتخاب منتظم: فأي طافر يفقد المؤثِّر أو يبدّله ينتشر في
المحصول كله في مواسم قليلة. والمربّون يكدّسون عدة مورّثات
مقاومة حتى يلزم فقد عدة مؤثِّرات دفعةً واحدة، ويخلطون
الأصناف، ويناوبون المورّثات على مدى سنين، ويحابون المقاومة
الواسعة المثارة بالأنماط التي لا تفلت منها طفرة واحدة.
\textbf{25.} $\varphi \approx 0.22$ تحت الظُلّة؛ \omterm{def:b3:plant-molecular-physiology:hormones}{والأوكسين} داخلًا :
خارجًا $\approx 43$؛ ونحو $170$ جزيء ماء مفقودًا لكل CO$_{2}$
مثبَّت.
\end{solution}
